 
\newtheorem{definition}{Definition}
\newtheorem{problem}{Problem}
\newtheorem{example}{Example}


%%%%%%%%%%% TEXT MACROS %%%%%%%%%%

\newcommand{\eg}{\textit{e.g.}}
\newcommand{\ie}{\textit{i.e.}}
\newcommand{\etc}{\textit{etc.}}
\newcommand{\wrt}{\textit{w.r.t.}}
\newcommand{\etal}{\textit{et al.}}

%%%%%%%%%%% MATHEMATICAL SYMBOLS %%%%%%%%%%

\usepackage{amsmath}
\usepackage{amssymb}
\newcommand{\randquant}[0]{\raisebox{0.75em}{\rotatebox{180}{\text{\sf R}}}}
\newcommand{\true}{\top}
\newcommand{\false}{\bot}

% -------------------------------------------------------------------------
% From https://tex.stackexchange.com/questions/43008/absolute-value-symbols
\DeclarePairedDelimiter\abs{\lvert}{\rvert}%
\DeclarePairedDelimiter\norm{\lVert}{\rVert}%

% Don't know if we need the following?
% Swap the definition of \abs* and \norm*, so that \abs
% and \norm resizes the size of the brackets, and the 
% starred version does not.
\makeatletter
\let\oldabs\abs
\def\abs{\@ifstar{\oldabs}{\oldabs*}}
%
\let\oldnorm\norm
\def\norm{\@ifstar{\oldnorm}{\oldnorm*}}
\makeatother


%%%%%%%%%%% PSEUDOCODE %%%%%%%%%%

\usepackage{algorithm}
\usepackage{algorithmicx}
\usepackage[noend]{algpseudocode}	% for hyper cross referencing (algpseudocode MUST BE LOADED LATER)
\algnewcommand\algorithmicinput{\textbf{INPUT:}}
\algnewcommand\INPUT{\item[\algorithmicinput]}

\algnewcommand\algorithmicoutput{\textbf{OUTPUT:}}
\algnewcommand\OUTPUT{\item[\algorithmicoutput]}

\algnewcommand{\IIf}[1]{\State\algorithmicif\ #1\ \algorithmicthen}
\algnewcommand{\EndIIf}{\unskip\ \algorithmicend\ \algorithmicif}
\algnewcommand{\IFor}[1]{\State\algorithmicfor\ #1\ \algorithmicdo}
\algnewcommand{\EndIFor}{\unskip\ \algorithmicend\ \algorithmicfor}
\algnewcommand{\IfThenElse}[3]{% \IfThenElse{<if>}{<then>}{<else>}
  \State \algorithmicif\ #1\ \algorithmicthen\ #2\ \algorithmicelse\ #3}

\algnewcommand\algorithmicswitch{\textbf{switch}}
\algnewcommand\algorithmiccase{\textbf{case}}

% New "environments"
\algdef{SE}[SWITCH]{Switch}{EndSwitch}[1]{\algorithmicswitch\ #1\ \algorithmicdo}{\algorithmicend\ \algorithmicswitch}%
\algdef{SE}[CASE]{Case}{EndCase}[1]{\algorithmiccase\ #1}{\algorithmicend\ \algorithmiccase}%
\algtext*{EndSwitch}%
\algtext*{EndCase}%

\makeatletter
\newcounter{algorithmicH}% New algorithmic-like hyperref counter
\let\oldalgorithmic\algorithmic
\renewcommand{\algorithmic}{%
  \stepcounter{algorithmicH}% Step counter
  \oldalgorithmic}% Do what was always done with algorithmic environment
\renewcommand{\theHALG@line}{ALG@line.\thealgorithmicH.\arabic{ALG@line}}
\makeatother


%%%%%%%%%%% SOLVERS AND LANGUAGES %%%%%%%%%%


\newcommand{\addmc}[0]{{\sf ADDMC}}
\newcommand{\approxcount}[0]{{\sf ApproxCount}}
\newcommand{\approxmc}[0]{{\sf ApproxMC4}}
\newcommand{\approxmcp}[0]{{\sf ApproxMC-p}}
\newcommand{\arjun}[0]{{\sf Arjun}}
\newcommand{\bpluse}[0]{{\sf B+E}}
\newcommand{\ctod}[0]{{\sf C2D}}
\newcommand{\cachet}[0]{{\sf cachet}}
\newcommand{\cachetwmc}[0]{{\sf cachet-wmc}}
\newcommand{\cnftoeadt}[0]{{\sf cnf2eadt}}
\newcommand{\cnftoobdd}[0]{{\sf CNF2OBDD}}
\newcommand{\cudd}[0]{{\sf CUDD}}
\newcommand{\cryptominisat}[0]{{\sf CryptoMiniSAT}}
\newcommand{\dfour}[0]{{\sf D4}}
\newcommand{\dpmc}[0]{{\sf DPMC}}
\newcommand{\dsharp}[0]{{\sf DSHARP}}
\newcommand{\gpmc}[0]{{\sf GPMC}}
\newcommand{\exactmc}[0]{{\sf ExactMC}}
\newcommand{\focusedflatsat}[0]{{\sf FocusedFlatSAT}}
\newcommand{\ganak}[0]{{\sf GANAK}}
\newcommand{\mbound}[0]{{\sf MBound}}
\newcommand{\minictod}[0]{{\sf miniC2D}}
\newcommand{\miniSAT}[0]{{\sf miniSAT}}
\newcommand{\python}[0]{{\sf Python}}
\newcommand{\samplecount}[0]{{\sf SampleCount}}
\newcommand{\sdd}[0]{{\sf SDD}}
\newcommand{\searchtreesampler}[0]{{\sf Search Tree Sampler}}
\newcommand{\sharpsat}[0]{{\sf SharpSAT}}
\newcommand{\sharpsattd}[0]{{\sf SharpSAT-TD}}
\newcommand{\treeofbdds}[0]{{\sf Tree-of-BDDs}}
\newcommand{\weightMC}[0]{{\sf WeightMC}}
\newcommand{\weightedsharpsat}[0]{{\sf weighted sharpSAT}}


%%%%%%%%%%% CNF ENCODINGS %%%%%%%%%%

\newcommand{\darohtwo}[0]{{\bf Dar02}}
\newcommand{\sbkohfive}[0]{{\bf SBK05}}
\newcommand{\cdohfive}[0]{{\bf CD05}}
\newcommand{\cdohsix}[0]{{\bf CD06}}
\newcommand{\bklmsixteen}[0]{{\bf BLKM16}}









